\section{Rezultati i rasprava}
\label{sec:rezultati}

Postupak je primijenjen na 300 deformacijskih stanja vlačnog pokusa. U nastavku su
prikazani rekonstruirano polje 3D pomaka, njegove komponente, evolucija pomaka kroz
pokus te polje površinskih deformacija, uz raspravu točnosti. Svi rezultati
uključuju čišćenje rubnih čvorova opisano u odjeljku \ref{ssec:ciscenje}.

\subsection{Polje trodimenzijskih pomaka}

Slika \ref{fig:dispmag} prikazuje deformirani oblik mjerne površine obojen iznosom
ukupnog 3D pomaka. Rekonstruirana površina jasno je zakrivljena, čime se potvrđuje da
mjerenje obuhvaća i izvanravninsku komponentu koju 2D korelacija ne bi mogla razlučiti.
Polje pomaka mijenja se glatko po površini, bez naglih diskontinuiteta, što upućuje na
konzistentnost rekonstrukcije; blaga neravnost uz sam rub zone odgovara preostaloj
(subpikselnoj) mjernoj nesigurnosti rubnih elemenata.

\begin{figure}[H]
  \centering
  \includegraphics[width=0.78\textwidth]{disp_magnitude.png}
  \caption{Deformirani 3D oblik mjerne površine obojen iznosom ukupnog pomaka
  $|\mathbf{U}|$ (mjerne jedinice kalibracije), za posljednje deformacijsko stanje.}
  \label{fig:dispmag}
\end{figure}

Razlaganje pomaka na komponente (Slika \ref{fig:dispcomp}) daje jasniju fizikalnu
sliku gibanja. Komponenta $U_x$ (uzduž kraće osi mjerne zone) praktički je
zanemariva. Ravninsko gibanje gotovo u cijelosti nosi komponenta $U_y$, pozitivna i
približno jednolika po površini, što odgovara vlačnom razvlačenju uzorka u smjeru
osi $y$ uz pomak mjerne zone kao cjeline. Izvanravninska komponenta $U_z$ posvuda je
negativna -- cijela se površina tijekom pokusa pomiče prema kamerama odnosno
,,spljošćuje``, s najvećim iznosom na zakrivljenijem dijelu zone. Upravo ta
kombinacija (usporedivi $U_y$ i $U_z$) čini uzorak reprezentativnim testom stereo
metode: jednokamerna 2D korelacija interpretirala bi izvanravninski pomak kao lažnu
ravninsku deformaciju.

\begin{figure}[H]
  \centering
  \includegraphics[width=\textwidth]{disp_components.png}
  \caption{Komponente 3D pomaka $U_x$, $U_y$ i $U_z$ za posljednje deformacijsko
  stanje, na zajedničkoj divergentnoj skali (kalibracijske jedinice). Ravninsko
  gibanje dominantno nosi $U_y$, a izvanravninsko $U_z$.}
  \label{fig:dispcomp}
\end{figure}

\subsection{Evolucija pomaka}

Slika \ref{fig:evolution} prikazuje najveći ukupni pomak $|\mathbf{U}|$ i najveću
izvanravninsku komponentu $|U_z|$ kroz niz frameova. Prvih tridesetak frameova čini
plato na razini $\approx 0{,}003$ kalibracijskih jedinica: uzorak još nije značajno
opterećen pa izmjerene vrijednosti predstavljaju efektivnu razinu šuma mjernog
lanca. Nakon platoa obje krivulje rastu monotono i približno linearno, uz blago
ubrzanje pri kraju pokusa (nagib u završnom dijelu $1{,}21$ puta veći od nagiba u
sredini pokusa). Vrijedi istaknuti da je prije čišćenja rubnih čvorova krivulja
pokazivala prividno naglo ubrzanje oko 240.\ framea -- ono se pokazalo artefaktom
dekorelacije čvorova 110 i 229 (Slika \ref{fig:detekcija}), a ne stvarnim ponašanjem
uzorka. U posljednjem stanju najveći ukupni pomak iznosi približno $0{,}09$, a
najveća izvanravninska komponenta $0{,}08$ (u mjernim jedinicama kalibracije) --
izvanravninski je pomak dakle gotovo jednak ravninskom, što potvrđuje nužnost stereo
mjerenja za ovakav uzorak.

\begin{figure}[H]
  \centering
  \includegraphics[width=0.80\textwidth]{disp_evolution.png}
  \caption{Evolucija najvećeg ukupnog pomaka $|\mathbf{U}|$ i najveće izvanravninske
  komponente $|U_z|$ kroz deformacijska stanja.}
  \label{fig:evolution}
\end{figure}

\subsection{Polje deformacija}

Slika \ref{fig:strain} prikazuje polje glavne deformacije $E_1$ po površini u
posljednjem stanju. Tipična (medijalna) vrijednost glavne deformacije iznosi približno
2\,\% ($1{,}95\,\%$), što je red veličine očekivan za vlačni pokus metalnog uzorka u
elastično-plastičnom području. Deformacija računata po elementu pokazuje razinu šuma
svojstvenu deriviranju polja pomaka na linearnim trokutastim elementima; skala boja
stoga je ograničena robusnim 97.\ percentilom ($12{,}04\,\%$). Nakon čišćenja rubnih
čvorova najveća vrijednost po elementu iznosi $27{,}36\,\%$ (prije čišćenja
nefizikalnih $328{,}69\,\%$); preostale povišene vrijednosti uz rub zone odraz su
subpikselne nesigurnosti rubnih elemenata koja se deriviranjem pojačava, a mogle bi
se dodatno smanjiti zaglađivanjem (ekstrapolacijom deformacija u čvorove) ili
profinjenjem mreže.

\begin{figure}[H]
  \centering
  \includegraphics[width=0.78\textwidth]{strain_E1.png}
  \caption{Polje glavne (najveće) deformacije $E_1$ na rekonstruiranoj 3D površini.}
  \label{fig:strain}
\end{figure}

\subsection{Procjena točnosti}
\label{ssec:tocnost}

Točnost sustava ocijenjena je posredno, kroz tri neovisna pokazatelja. Prvo,
rekonstrukcija referentnog oblika podudara se sa zadanom mrežom do strojne
preciznosti ($\sim 10^{-15}$), čime je isključena sustavna pogreška triangulacije.
Drugo, srednja reprojekcijska pogreška trianguliranih deformiranih točaka iznosi
približno $0{,}15$ piksela (medijan $0{,}13$ px, najveća $1{,}10$ px) -- budući da ta
veličina objedinjuje pogrešku korelacije, kalibracije i triangulacije, subpikselna
vrijednost potvrđuje visoku ukupnu točnost mjerenja i u razini je vrijednosti koje
se u literaturi navode za dobro kalibrirane stereo-DIC sustave. Treće, plato
krivulje evolucije u neopterećenom dijelu pokusa daje procjenu šuma pomaka od
$\approx 0{,}003$ kalibracijskih jedinica.

Lokalizirani gubitak korelacije na tri rubna čvora kamere 2, koji je prije čišćenja
bio dominantan izvor pogreške (reprojekcija do $6{,}25$ px, nefizikalni ekstremi
deformacija), automatski je detektiran i saniran postupkom iz odjeljka
\ref{ssec:ciscenje}. Preostala nesigurnost svodi se na šum polja deformacija po
elementu, koji se po potrebi može dodatno smanjiti zaglađivanjem ili profinjenjem
mreže uz rub.
